% ---
% titre: Calcul flash: 9VP - cycle 1 (semaines 1-6) - théorie
% theme: NO
% chapitre: Nombres naturels et décimaux
% sous_chapitre: Calcul réfléchi et mental
% niveau: [9e]
% type: TH
% tags: [c-division,c-livrets, c-multiplication,c-operations,c-priorite-operations,c-puissance, f-individuel,p-flash-calcul]
% difficulte: 1
% source: latex
% ---

\newtcolorbox{regle}[1]{enhanced, sharp corners, boxrule=0pt, interior hidden, frame hidden,borderline west={2.2pt}{0pt}{themeColor},left=3mm, right=2mm, top=2mm, bottom=2mm,title={\textbf{\textcolor{themeColor}{#1}}}, fonttitle=\bfseries,attach title to upper={\\},}

\newtcolorbox{astuce}[1]{enhanced, sharp corners, halign=justify, valign=top, colback=white, colframe=themeColor, boxrule=2.2pt, left=3mm, right=2mm, top=2mm, bottom=2mm,title={\textbf{\textcolor{themeColor}{#1}}}, fonttitle=\bfseries\small,attach title to upper={\\},}

\newtcolorbox{attention}{enhanced, sharp corners, boxrule=0pt,colback=red!5!white, colframe=red!5!white,borderline west={2.2pt}{0pt}{red!70!black},left=3mm, right=2mm, top=2mm, bottom=2mm,}

\newpage
\header{Calcul mental: 9VP - Semaine 1 - Tables \& livrets} 

 \vfill
%-----------------------------------------------------------------------
\subsection*{1. Les tables de multiplication}
 
\begin{regle}{Définition}
La table de $n$ liste tous les multiples de $n$ : $n \cdot 1, n \cdot 2, \ldots, n \cdot 12$.
\end{regle}
 
\vspace{6pt}
 \textbf{Tables à connaître par coeur}:
 
\vspace{4pt}
\def\arraystretch{1.6}
\begin{tabularx}{\linewidth}{|>{\centering\arraybackslash}p{1.4cm}|*{12}{>{\centering\arraybackslash}X|}}
\hline
\cellcolor{themeNO!50}$\cdot$ &\cellcolor{themeNO!30}1 & \cellcolor{themeNO!30}2 & \cellcolor{themeNO!30}3 & \cellcolor{themeNO!30}4 &\cellcolor{themeNO!30}5 & \cellcolor{themeNO!30}6 & \cellcolor{themeNO!30}7 & \cellcolor{themeNO!30}8 &\cellcolor{themeNO!30}9 & \cellcolor{themeNO!30}10 & \cellcolor{themeNO!30}11 & \cellcolor{themeNO!30}12 \\
\hline
\cellcolor{themeNO!30}1  & 1 & 2 & 3 & 4 & 5 & 6 & 7 & 8 & 9 & 10 & 11 & 12 \\\hline
\cellcolor{themeNO!30}2  & 2 & 4 & 6 & 8 & 10 & 12 & 14 & 16 & 18 & 20 & 22 & 24 \\\hline
\cellcolor{themeNO!30}3  & 3 & 6 & 9 & 12 & 15 & 18 & 21 & 24 & 27 & 30 & 33 & 36 \\\hline
\cellcolor{themeNO!30}4  & 4 & 8 & 12 & 16 & 20 & 24 & 28 & 32 & 36 & 40 & 44 & 48 \\\hline
\cellcolor{themeNO!30}5  & 5 & 10 & 15 & 20 & 25 & 30 & 35 & 40 & 45 & 50 & 55 & 60 \\\hline
\cellcolor{themeNO!30}6  & 6 & 12 & 18 & 24 & 30 & 36 & 42 & 48 & 54 & 60 & 66 & 72 \\\hline
\cellcolor{themeNO!30}7  & 7 & 14 & 21 & 28 & 35 & 42 & 49 & 56 & 63 & 70 & 77 & 84 \\\hline
\cellcolor{themeNO!30}8  & 8 & 16 & 24 & 32 & 40 & 48 & 56 & 64 & 72 & 80 & 88 & 96 \\\hline
\cellcolor{themeNO!30}9  & 9 & 18 & 27 & 36 & 45 & 54 & 63 & 72 & 81 & 90 & 99 & 108 \\\hline
\cellcolor{themeNO!30}10  & 10 & 20 & 30 & 40 & 50 & 60 & 70 & 80 & 90 & 100 & 110 & 120 \\\hline
\cellcolor{themeNO!30}11 & 11 & 22 & 33 & 44 & 55 & 66 & 77 & 88 & 99 & 110 & 121 & 132 \\\hline
\cellcolor{themeNO!30}12 & 12 & 24 & 36 & 48 & 60 & 72 & 84 & 96 & 108 & 120 & 132 & 144 \\\hline
\end{tabularx}
\def\arraystretch{1}

%-----------------------------------------------------------------------
\vspace{10pt}
\subsection*{2. Division}
 
\begin{regle}{Définition}
Si $a \cdot b = c$, alors $c \div a = b$ et $c \div b = a$.
\end{regle}
 
\vspace{6pt}
 
\begin{astuce}{Stratégie}
Pour calculer $63 \div 9$, se demander : \textit{\og $9$ fois combien font $63$ ?\fg}\\
On cherche dans la table de $9$ : $9 \cdot 7 = 63$, donc $63 \div 9 = 7$.
\end{astuce}
 
\vspace{8pt}
 
\newpage
 %-----------------------------------------------------------------------
\subsection*{3. Les carrés d'un nombre}
 
\begin{regle}{Définition}
Le carré d'un nombre $n$ est le produit de $n$ par lui-même : $n^2 = n \cdot n$ \quad (lire : $n$ au carré)
\end{regle}
 
\vspace{6pt}
 \textbf{Carrés à connaître par coeur}:
 
\vspace{4pt} 
\def\arraystretch{1.5}
\begin{tabularx}{0.3\linewidth}{|>{\centering\arraybackslash}X|>{\centering\arraybackslash}X|}
\hline
\cellcolor{themeNO!30}$n$ & \cellcolor{themeNO!30}$n^2$ \\\hline
$1$ & $1$  \\\hline
$2$ & $4$  \\\hline
$3$ & $9$  \\\hline
$4$ & $16$  \\\hline
$5$ & $25$ \\\hline
$6$ & $36$ \\\hline
\end{tabularx}
\hfill
\begin{tabularx}{0.3\linewidth}{|>{\centering\arraybackslash}X|>{\centering\arraybackslash}X|}
\hline
\cellcolor{themeNO!30}$n$ & \cellcolor{themeNO!30}$n^2$ \\\hline
$7$ & $49$ \\\hline
$8$ & $64$ \\\hline
$9$ & $81$ \\\hline
$10$ & $100$ \\\hline
$11$ & $121$ \\\hline
$12$ & $144$  \\\hline
\end{tabularx}
\hfill
\begin{tabularx}{0.3\linewidth}{|>{\centering\arraybackslash}X|>{\centering\arraybackslash}X|}
\hline
\cellcolor{themeNO!30}$n$ & \cellcolor{themeNO!30}$n^2$ \\\hline
$13$ & $169$ \\\hline
$14$ & $196$ \\\hline
$15$ & $225$ \\\hline
$20$ & $400$ \\\hline
$100$ & $10\,000$ \\\hline
$1000$ & $1\,000\,000$ \\\hline
\end{tabularx}
\def\arraystretch{1}
   
%-----------------------------------------------------------------------
\vspace{10pt}
\subsection*{4. Astuces pour retenir les tables}
 
 \begin{tcbraster}[raster columns=2, raster equal height, raster valign=top, raster column skip=4mm, raster row skip=4mm]
\begin{astuce}{Commutativité : $a \times b = b \times a$}
On peut toujours retourner une multiplication.\\[4pt]
\textit{Exemple :} Si on oublie $8 \times 7$, on calcule $7 \times 8 = 56$.
\end{astuce}
 \begin{astuce}{Table de 11 : le truc du miroir}
Pour $11 \times n$ (avec $n \leq 9$) : répéter le chiffre.\\[4pt]
\textit{Exemple :} $11 \times 7 = 77$\\
\phantom{\textit{Exemple :}} $11 \times 8 = 88$\\[6pt]
Pour $n > 9$ : $11 \times n = 10n + n$\\[4pt]
\textit{Exemple :} $11 \times 12 = 120 + 12 = 132$
\end{astuce}
 \begin{astuce}{Table de 9 : soustraire une fois}
$9 \times n = 10 \times n - n$\\[4pt]
\textit{Démonstration :} $9 = 10 - 1$, donc $9 \times n = (10-1) \times n = 10n - n$\\[2pt]
\textit{Exemple :} $9 \times 7 = 70 - 7 = 63$\\
\phantom{\textit{Exemple :}} $9 \times 8 = 80 - 8 = 72$
\end{astuce}
 \begin{astuce}{Table de 12 : ajouter une fois}
$12 \times n = 10 \times n + 2 \times n$\\[4pt]
\textit{Démonstration :} $12 = 10 + 2$, donc $12 \times n = 10n + 2n$\\[2pt]
\textit{Exemple :} $12 \times 8 = 80 + 16 = 96$\\
\phantom{\textit{Exemple :}} $12 \times 7 = 70 + 14 = 84$
\end{astuce}
 \end{tcbraster}
 
\vspace{8pt}
 
\begin{attention}
\textbf{Les tables difficiles à retenir} - entraîne-toi particulièrement sur :\\[4pt]
$6 \times 7 = 42$ \qquad $6 \times 8 = 48$ \qquad $7 \times 8 = 56$ \qquad $7 \times 9 = 63$
\end{attention}


\newpage
\header{Calcul mental: 9VP - Semaine 2 - $\times$ et $\div$ par 10, 100, 1000} 

%-----------------------------------------------------------------------
\subsection*{1. La règle des rangs}
 
\begin{regle}{Règle fondamentale}
Multiplier ou diviser par 10, 100 ou 1000 revient à \textbf{déplacer la virgule} d'un certain nombre de rangs.\\[6pt]
\begin{tabular}{lll}
$\times\, 10$ & $\to$ & virgule avance d'\textbf{1 rang} vers la droite \\
$\times\, 100$ & $\to$ & virgule avance de \textbf{2 rangs} vers la droite \\
$\times\, 1000$ & $\to$ & virgule avance de \textbf{3 rangs} vers la droite \\[4pt]
$\div\, 10$ & $\to$ & virgule recule d'\textbf{1 rang} vers la gauche \\
$\div\, 100$ & $\to$ & virgule recule de \textbf{2 rangs} vers la gauche \\
$\div\, 1000$ & $\to$ & virgule recule de \textbf{3 rangs} vers la gauche \\
\end{tabular}\\[6pt]
\textit{Astuce :} le nombre de \textbf{zéros} = le nombre de \textbf{rangs} à déplacer.
\end{regle}
 
\vspace{10pt}
 
%-----------------------------------------------------------------------
\subsection*{2. Exemples}
 
\def\arraystretch{1.7}
\begin{tabularx}{\linewidth}{|X|X|X|}
\hline
\cellcolor{themeNO!30}\textbf{Opération} & \cellcolor{themeNO!30}\textbf{Déplacement} & \cellcolor{themeNO!30}\textbf{Résultat} \\\hline
$3{,}7 \times 10$ & 1 rang à droite & $37$ \\\hline
$3{,}7 \times 100$ & 2 rangs à droite & $370$ \\\hline
$3{,}7 \times 1000$ & 3 rangs à droite & $3700$ \\\hline
$3{,}7 \div 10$ & 1 rang à gauche & $0{,}37$ \\\hline
$3{,}7 \div 100$ & 2 rangs à gauche & $0{,}037$ \\\hline
$0{,}08 \times 1000$ & 3 rangs à droite & $80$ \\\hline
$340 \div 1000$ & 3 rangs à gauche & $0{,}34$ \\\hline
\end{tabularx}
\def\arraystretch{1}
 
\vspace{10pt}
 
%-----------------------------------------------------------------------
\subsection*{3. Cas particulier : les entiers}
 
\begin{astuce}{Multiplier un entier par 10, 100, 1000}
Quand on multiplie un nombre \textbf{entier} par 10, 100 ou 1000, on \textbf{ajoute des zéros} à droite.\\[4pt]
\textit{Pourquoi ?} Un entier a sa virgule après le dernier chiffre : $7 = 7{,}0$\\
Déplacer de 1 rang à droite : $7{,}0 \to 70{,} = 70$\\[4pt]
$7 \times 10 = 70$ \qquad $7 \times 100 = 700$ \qquad $7 \times 1000 = 7000$
\end{astuce}
 
\vspace{8pt}
 
%-----------------------------------------------------------------------
\subsection*{4. Vérification}
 
\begin{regle}{Comment vérifier son résultat}
\begin{itemize}[leftmargin=*, itemsep=3pt, topsep=2pt, label=$\triangleright$]
  \item \textbf{Multiplier} $\to$ le résultat doit être \textbf{plus grand} que le nombre de départ.
  \item \textbf{Diviser} $\to$ le résultat doit être \textbf{plus petit} que le nombre de départ.
  \item Compter le nombre de zéros de 10, 100 ou 1000 et vérifier que la virgule a bien bougé d'autant de rangs.
\end{itemize}
\end{regle}
 
\vspace{8pt}
 
\begin{attention}
\textbf{Erreur fréquente :} $0{,}07 \times 100 \neq 0{,}7$\\
Compter les rangs : 2 rangs à droite depuis $0{,}07$ donne $07{,}$ = $7$. \\
\textbf{Résultat correct : $0{,}07 \times 100 = 7$}
\end{attention}

\newpage
\header{Calcul mental: 9VP - Semaine 3 - Doubles \& moitiés}
 
%-----------------------------------------------------------------------
\subsection*{1. Doubles et moitiés}
 
\begin{regle}{Définitions}
\begin{tabular}{lll}
\textbf{Doubler} & = & multiplier par $2$ \quad ($\times 2$) \\
\textbf{Prendre la moitié} & = & diviser par $2$ \quad ($\div 2$)\\
\end{tabular}
\end{regle}
 
\vspace{6pt}
 
\begin{tcbraster}[raster columns=2, raster equal height, raster valign=top,
  raster column skip=4mm, raster row skip=4mm]
 
\begin{astuce}{Doubler un décimal}
Doubler chaque partie séparément.\\[4pt]
$2 \times 1{,}4 = 2 \times 1 + 2 \times 0{,}4 = 2 + 0{,}8 = 2{,}8$\\[2pt]
$2 \times 3{,}5 = 6 + 1{,}0 = 7$
\end{astuce}
 
\begin{astuce}{Moitié d'un nombre impair}
La moitié d'un impair donne un nombre avec $,5$.\\[4pt]
$13 \div 2 = 6{,}5$ \qquad $7 \div 2 = 3{,}5$\\[2pt]
\textit{Astuce :} $13 = 12 + 1 \to 6 + 0{,}5 = 6{,}5$
\end{astuce}
 
\end{tcbraster}
 
\vspace{8pt}
 
%-----------------------------------------------------------------------
\subsection*{2. Liens avec $\times\, 0{,}5$ et $\div\, 0{,}5$}
 
\begin{regle}{Équivalences importantes}
\begin{tabular}{lllll}
$\times\, 0{,}5$ & $=$ & $\div 2$ & $=$ & prendre la moitié \\[4pt]
$\div\, 0{,}5$ & $=$ & $\times 2$ & $=$ & doubler \\
\end{tabular}\\[6pt]
\textit{Pourquoi ?} $0{,}5 = \dfrac{1}{2}$, donc multiplier par $0{,}5$ revient à multiplier par $\dfrac{1}{2}$, c'est-à-dire diviser par $2$.\\[4pt]
Et diviser par $\dfrac{1}{2}$ revient à multiplier par $2$.
\end{regle}
 
\vspace{8pt}
 
\def\arraystretch{1.6}
\begin{tabularx}{\linewidth}{|X|X|}
\hline
\cellcolor{themeNO!30}\textbf{Expression} & \cellcolor{themeNO!30}\textbf{Calcul} \\\hline
$0{,}5 \times 48$ & $48 \div 2 = 24$ \\\hline
$0{,}5 \times 130$ & $130 \div 2 = 65$ \\\hline
$124 \div 0{,}5$ & $124 \times 2 = 248$ \\\hline
$30 \div 0{,}5$ & $30 \times 2 = 60$ \\\hline
\end{tabularx}
\def\arraystretch{1}
 
\vspace{10pt}
 
%-----------------------------------------------------------------------
\subsection*{3. Enchaîner les doubles : $\times 4$}
 
\begin{astuce}{$\times\, 4$ = doubler deux fois}
\textit{Démonstration :} $4 = 2 \times 2$, donc $n \times 4 = (n \times 2) \times 2$\\[4pt]
\textit{Exemple :} $13 \times 4 \to 13 \times 2 = 26 \to 26 \times 2 = 52$\\[2pt]
De même $\div 4$ = prendre la moitié deux fois.
\end{astuce}
 
\vspace{8pt}
 
\begin{attention}
\textbf{Ne pas confondre :}\\[4pt]
$\times\, 0{,}5 = \div 2$ \quad (résultat \textbf{plus petit})\\
$\div\, 0{,}5 = \times 2$ \quad (résultat \textbf{plus grand})\\[4pt]
Vérifier : diviser par un nombre inférieur à 1 donne un résultat \textbf{plus grand} !
\end{attention}

\newpage
\header{Calcul mental: 9VP - Semaine 4 - Astuces $\times$5, $\times$9, $\times$11, $\times$25}

\subsection*{1. Multiplier et diviser par 5}

\begin{regle}{Les astuces pour $\times 5$ et $\div 5$}
\begin{tabular}{lllll}
$n \times 5$ & $=$ & $n \times 10 \div 2$ & $=$ & multiplier par 10, puis prendre la moitié \\[4pt]
$n \div 5$ & $=$ & $n \times 2 \div 10$ & $=$ & doubler, puis diviser par 10 \\
\end{tabular}\\[6pt]
\textit{Pourquoi ?} $5 = \dfrac{10}{2}$, donc $n \times 5 = n \times \dfrac{10}{2} = \dfrac{n \times 10}{2}$
\end{regle}

\vspace{6pt}

\def\arraystretch{1.6}
\begin{tabularx}{\linewidth}{|X|X|X|}
\hline
\cellcolor{themeNO!30}\textbf{Calcul} & \cellcolor{themeNO!30}\textbf{Étapes} & \cellcolor{themeNO!30}\textbf{Résultat} \\\hline
$34 \times 5$ & $34 \times 10 = 340 \quad \to \quad 340 \div 2$ & $170$ \\\hline
$22 \times 5$ & $22 \times 10 = 220 \quad \to \quad 220 \div 2$ & $110$ \\\hline
$46 \div 5$ & $46 \times 2 = 92 \quad \to \quad 92 \div 10$ & $9{,}2$ \\\hline
$70 \div 5$ & $70 \times 2 = 140 \quad \to \quad 140 \div 10$ & $14$ \\\hline
\end{tabularx}
\def\arraystretch{1}

\vspace{10pt}

%-----------------------------------------------------------------------
\subsection*{2. Multiplier par 9 et par 11}

\begin{tcbraster}[raster columns=2, raster equal height, raster valign=top,
  raster column skip=4mm, raster row skip=4mm]

\begin{regle}{$\times\, 9$ : multiplier par 10, puis soustraire}
$n \times 9 = n \times 10 - n$\\[4pt]
\textit{Pourquoi ?} $9 = 10 - 1$, donc $n \times 9 = n \times (10-1) = 10n - n$\\[6pt]
$7 \times 9 = 70 - 7 = \mathbf{63}$\\
$18 \times 9 = 180 - 18 = \mathbf{162}$\\
$25 \times 9 = 250 - 25 = \mathbf{225}$
\end{regle}

\begin{regle}{$\times\, 11$ : multiplier par 10, puis ajouter}
$n \times 11 = n \times 10 + n$\\[4pt]
\textit{Pourquoi ?} $11 = 10 + 1$, donc $n \times 11 = n \times (10+1) = 10n + n$\\[6pt]
$7 \times 11 = 70 + 7 = \mathbf{77}$\\
$13 \times 11 = 130 + 13 = \mathbf{143}$\\
$23 \times 11 = 230 + 23 = \mathbf{253}$
\end{regle}

\end{tcbraster}

\vspace{10pt}

%-----------------------------------------------------------------------
\subsection*{3. Multiplier par 25}

\begin{regle}{$\times\, 25$ : multiplier par 100, puis diviser par 4}
$n \times 25 = n \times 100 \div 4$\\[4pt]
\textit{Pourquoi ?} $25 = \dfrac{100}{4}$, donc $n \times 25 = n \times \dfrac{100}{4} = \dfrac{n \times 100}{4}$\\[4pt]
\textit{Astuce :} diviser par 4 = prendre la moitié deux fois.
\end{regle}

\vspace{6pt}

\def\arraystretch{1.6}
\begin{tabularx}{\linewidth}{|X|X|X|}
\hline
\cellcolor{themeNO!30}\textbf{Calcul} & \cellcolor{themeNO!30}\textbf{Étapes} & \cellcolor{themeNO!30}\textbf{Résultat} \\\hline
$24 \times 25$ & $24 \times 100 = 2400 \to 2400 \div 4$ & $600$ \\\hline
$36 \times 25$ & $36 \times 100 = 3600 \to 3600 \div 4$ & $900$ \\\hline
$16 \times 25$ & $16 \times 100 = 1600 \to 1600 \div 4$ & $400$ \\\hline
$48 \times 25$ & $48 \times 100 = 4800 \to 4800 \div 4$ & $1200$ \\\hline
\end{tabularx}
\def\arraystretch{1}

\vspace{10pt}

%-----------------------------------------------------------------------
\subsection*{4. Récapitulatif des astuces}

\def\arraystretch{1.7}
\begin{tabularx}{\linewidth}{|>{\centering\arraybackslash}p{2cm}|X|X|}
\hline
\cellcolor{themeNO!30}\textbf{Opération} & \cellcolor{themeNO!30}\textbf{Astuce} & \cellcolor{themeNO!30}\textbf{Exemple rapide} \\\hline
$\times\, 5$ & $\times 10$ puis $\div 2$ & $34 \times 5 : 340 \div 2 = 170$ \\\hline
$\div\, 5$ & $\times 2$ puis $\div 10$ & $46 \div 5 : 92 \div 10 = 9{,}2$ \\\hline
$\times\, 9$ & $\times 10$ puis $- n$ & $8 \times 9 : 80 - 8 = 72$ \\\hline
$\times\, 11$ & $\times 10$ puis $+ n$ & $8 \times 11 : 80 + 8 = 88$ \\\hline
$\times\, 25$ & $\times 100$ puis $\div 4$ & $32 \times 25 : 3200 \div 4 = 800$ \\\hline
\end{tabularx}
\def\arraystretch{1}

%%%%%
\newpage
\header{Calcul mental: 9VP - Semaine 5 - Priorité des opération}
 

%-----------------------------------------------------------------------
\subsection*{1. La règle de priorité}
 
\begin{regle}{Ordre des opérations (sans parenthèses)}
Dans une expression sans parenthèses, on calcule \textbf{toujours dans cet ordre} :\\[6pt]
\begin{center}
\large $\times$ et $\div$ \textbf{avant} $+$ et $-$
\end{center}
\vspace{4pt}
Si plusieurs $\times$ et $\div$ se suivent : on calcule \textbf{de gauche à droite}.\\
Si plusieurs $+$ et $-$ se suivent : on calcule \textbf{de gauche à droite}.
\end{regle}
 
\vspace{8pt}
 
%-----------------------------------------------------------------------
\subsection*{2. Exemples détaillés}
 
\def\arraystretch{1.8}
\begin{tabularx}{\linewidth}{|X|X|}
\hline
\cellcolor{themeNO!30}\textbf{Expression} & \cellcolor{themeNO!30}\textbf{Calcul étape par étape} \\\hline
$3 + 4 \times 2$ &
  $= 3 + \underbrace{4 \times 2}_{8} = 3 + 8 = \mathbf{11}$ \\\hline
$10 - 6 \div 2$ &
  $= 10 - \underbrace{6 \div 2}_{3} = 10 - 3 = \mathbf{7}$ \\\hline
$7 \times 4 - 5 \times 3$ &
  $= \underbrace{7 \times 4}_{28} - \underbrace{5 \times 3}_{15} = 28 - 15 = \mathbf{13}$ \\\hline
$18 \div 6 + 4 \times 2$ &
  $= \underbrace{18 \div 6}_{3} + \underbrace{4 \times 2}_{8} = 3 + 8 = \mathbf{11}$ \\\hline
$12 \div 4 \times 3$ &
  $= \underbrace{12 \div 4}_{3} \times 3 = 3 \times 3 = \mathbf{9}$ \quad (gauche à droite)\\\hline
\end{tabularx}
\def\arraystretch{1}
 
\vspace{10pt}
 
%-----------------------------------------------------------------------
\subsection*{3. Stratégie de calcul}
 
\begin{astuce}{Méthode en 2 temps}
\textbf{Étape 1 :} Repérer (ou encercler) tous les $\times$ et $\div$ dans l'expression.\\
\textbf{Étape 2 :} Les calculer en premier, puis finir avec les $+$ et $-$.\\[6pt]
\textit{Exemple :} $\quad 5 + \boxed{8 \times 2} - 1 = 5 + 16 - 1 = 20$
\end{astuce}
 
\vspace{8pt}
 
%-----------------------------------------------------------------------
\subsection*{4. Erreurs fréquentes}
 
\begin{attention}
\textbf{Erreur :} calculer de gauche à droite sans tenir compte des priorités.\\[4pt]
$3 + 4 \times 2 \neq 7 \times 2 = 14$ \quad {\small (FAUX)}\\
$3 + 4 \times 2 = 3 + 8 = 11$ \quad {\small (CORRECT)}\\[6pt]
\textbf{Erreur :} oublier que $\div$ a la même priorité que $\times$.\\[4pt]
$12 \div 4 \times 3 \neq 12 \div 12 = 1$ \quad {\small (FAUX — on ne peut pas regrouper $4 \times 3$)}\\
$12 \div 4 \times 3 = 3 \times 3 = 9$ \quad {\small (CORRECT — de gauche à droite)}
\end{attention}
 
\vspace{10pt}
 
%-----------------------------------------------------------------------
\subsection*{5. Tableau de synthèse}
 
\def\arraystretch{1.6}
\begin{tabularx}{\linewidth}{|>{\centering\arraybackslash}p{3.5cm}|>{\centering\arraybackslash}p{3.5cm}|X|}
\hline
\cellcolor{themeNO!30}\textbf{Priorité} & \cellcolor{themeNO!30}\textbf{Opérations} & \cellcolor{themeNO!30}\textbf{Ordre} \\\hline
1\textsuperscript{re} (plus haute) & $\times \quad \div$ & De gauche à droite \\\hline
2\textsuperscript{e} (plus basse) & $+ \quad -$ & De gauche à droite \\\hline
\end{tabularx}
\def\arraystretch{1}